% Cooled-BGO CNN position resolution, prepared as a paper fragment.
%
% Required packages in the parent document: graphicx, booktabs, amsmath, url.
%
% Usage:
%   % Cooled-BGO CNN position resolution, prepared as a paper fragment.
%
% Required packages in the parent document: graphicx, booktabs, amsmath, url.
%
% Usage:
%   % Cooled-BGO CNN position resolution, prepared as a paper fragment.
%
% Required packages in the parent document: graphicx, booktabs, amsmath, url.
%
% Usage:
%   % Cooled-BGO CNN position resolution, prepared as a paper fragment.
%
% Required packages in the parent document: graphicx, booktabs, amsmath, url.
%
% Usage:
%   \input{cnn_bgo}             % inserts the main-text section below
%   ...
%   \appendix
%   \cnnbgoappendix             % inserts the methods appendix
%
% The figure directory can be overridden before this file is input:
%   \renewcommand{\cnnbgofigdir}{path/to/cnn_bgo_win}

\providecommand{\cnnbgofigdir}{../analysis/results/cnn/cnn_bgo_win}
\providecommand{\cnncode}[1]{\texttt{#1}}
\providecommand{\cnnpath}[1]{\path{#1}}

% =====================================================================
% Main-text section

\section{Position resolution in cooled BGO with a convolutional network}
\label{sec:cnn-bgo}

The three-dimensional interaction position was reconstructed from the
$8\times8$ SiPM light map with CryspNet, a compact convolutional neural network.
The result reported here uses the complete held-out BGO sample in the standard
511-keV analysis window. ``No selection'' means that every event in this window
is retained: no reconstructability or CNN-quality requirement is applied. The
test sample contains $54\,021$ events and includes the photoelectric and Compton
topologies present in the window.

The residual distributions have a narrow central population and a broader tail,
so a single RMS does not fully characterize the response. The central 68\%
absolute-residual quantiles are $1.08$, $1.08$ and $1.26$~mm in $x$, $y$ and
$z$, respectively; the corresponding global RMS values are $1.62$, $1.62$ and
$2.46$~mm. In three dimensions, $9.73\%$ of the reconstructed points lie more
than 5~mm from the true first interaction.

A compact detector-response description is obtained by fitting each coordinate
with two Gaussian components sharing a mean,
\begin{equation}
  p(\Delta q) = f_{\rm c}\,
  \mathcal{N}(\mu,\sigma_{\rm c})
  + (1-f_{\rm c})\,
  \mathcal{N}(\mu,\sigma_{\rm t}),
  \qquad q\in\{x,y,z\}.
  \label{eq:cnn-bgo-two-gauss}
\end{equation}
Table~\ref{tab:cnn-bgo-resolution} gives the complete no-selection result. The
Gaussian core has widths $0.80$, $0.82$ and $0.96$~mm, corresponding to FWHM
values $1.88$, $1.92$ and $2.26$~mm. It contains $72$--$82\%$ of the events,
while the remaining response is described by a component with
$\sigma_{\rm t}\simeq2.5$--$2.6$~mm. Thus cooled BGO provides a sub-millimetre
Gaussian core in all three coordinates and a typical-event resolution of about
1.1--1.3~mm when the full, unselected window sample is considered.

\begin{table}[htbp]
  \centering
  \begin{tabular}{lrrrrrr}
    \toprule
    coordinate
      & $p_{68}$ & RMS
      & $\sigma_{\rm c}$ & FWHM$_{\rm c}$
      & $f_{\rm c}$ & $\sigma_{\rm t}$ \\
      & \multicolumn{2}{c}{[mm]}
      & \multicolumn{2}{c}{[mm]}
      & & [mm] \\
    \midrule
    $x$ & 1.083 & 1.619 & 0.796 & 1.875 & 0.716 & 2.505 \\
    $y$ & 1.081 & 1.622 & 0.816 & 1.922 & 0.737 & 2.586 \\
    $z$ & 1.256 & 2.464 & 0.959 & 2.258 & 0.821 & 2.588 \\
    \bottomrule
  \end{tabular}
  \caption{No-selection CNN position response in BGO. The first two columns
  describe the full residual distribution; the remaining columns are the
  two-Gaussian parameters of Eq.~\eqref{eq:cnn-bgo-two-gauss}.}
  \label{tab:cnn-bgo-resolution}
\end{table}

\begin{figure}[htbp]
  \centering
  \includegraphics[width=\textwidth]{\cnnbgofigdir/residual_fit.png}
  \caption{CNN residuals for the complete no-selection BGO test sample in
  $x$, $y$ and $z$ (steps). The solid curve is the two-component Gaussian fit;
  the dashed and dotted curves show its core and tail terms. The logarithmic
  ordinate exposes the broad response beneath the sub-millimetre core. The
  modest positive skew in the depth residual is retained in the data and is
  summarized, rather than reproduced in detail, by the symmetric fit.}
  \label{fig:cnn-bgo-residual-fit}
\end{figure}

% =====================================================================
% Appendix block. Invoke with \cnnbgoappendix after the parent paper's
% \appendix command.

\newcommand{\cnnbgoappendix}{%
\section{CNN reconstruction and response fit}
\label{app:cnn-bgo}

\subsection{Network and inputs}

Each event is represented by its $8\times8$ photoelectron map and its total
photoelectron count. The map is divided by its sum, separating the spatial light
pattern from the event amplitude. The standardized logarithm of the total count
is supplied as a second, scalar input. The three target coordinates are mapped
linearly to $[-1,1]$ during training and converted back to millimetres for the
reported residuals.

CryspNet begins with a $3\times3$ convolution producing 64 channels, followed by
two 64-channel residual blocks at the native $8\times8$ resolution. A stride-two
convolution increases the representation to 128 channels and reduces it to
$4\times4$; two further residual blocks are followed by global average pooling.
The pooled 128-component descriptor is concatenated with the scalar light-yield
input and passed through a 128-unit fully connected layer and a linear
three-coordinate output. Batch normalization and ReLU activations are used
throughout the convolutional path. The network contains $830\,531$ trainable
parameters.

\subsection{Training outcome}

The window sample contains $360\,135$ events. A seeded permutation divides it
into 70\% training, 15\% validation and 15\% test subsets; the test subset is
held out until the final evaluation. The square detector symmetry supplies the
eight exact rotations and reflections of every training map, with the same
transformations applied to the transverse target coordinates. Validation and
test events remain unaugmented.

Training minimizes a Huber loss on the normalized coordinates with AdamW,
using a batch size of 512, an initial learning rate of $10^{-3}$, weight decay
$10^{-4}$ and a cosine learning-rate schedule over 15 epochs. The training loss
falls from $0.0500$ to $0.0390$. The saved checkpoint is the epoch-10 state,
which minimizes the mean validation RMSE; its coordinate validation values are
$1.608$, $1.612$ and $2.434$~mm. The independent test result is the response
reported in Table~\ref{tab:cnn-bgo-resolution}.

\subsection{Two-component response fit}

The fit uses the CNN prediction minus the true first-interaction coordinate for
each of the $54\,021$ test events. Each residual is histogrammed in 150 bins over
$[-15,15]$~mm and fitted with Eq.~\eqref{eq:cnn-bgo-two-gauss}, using
Poisson weights. The two components share the fitted mean; the narrower width is
identified as the core. The fitted means are $-0.032$, $-0.058$ and
$-0.127$~mm for $x$, $y$ and $z$, showing a small global bias. The depth
distribution also contains a positive, non-Gaussian shoulder from multi-site
events whose light is dominated by a deeper interaction. Equation
\eqref{eq:cnn-bgo-two-gauss} supplies a compact core--tail parameterization,
while Fig.~\ref{fig:cnn-bgo-residual-fit} displays the full residual shape.

The stored configuration is
\cnnpath{analysis/cnn/runs/cnn_bgo_win.toml}. The training history, metrics,
held-out events and fit products are stored under
\cnnpath{analysis/results/cnn/cnn_bgo_win/} as \cnnpath{history.json},
\cnnpath{metrics.json}, \cnnpath{test_events.h5}, \cnnpath{residual_fit.json}
and \cnnpath{residual_fit.png}. The response fit is produced from the held-out
event table by \cnnpath{analysis/cnn/fit_residuals.py}; it does not retrain the
network.
}
             % inserts the main-text section below
%   ...
%   \appendix
%   \cnnbgoappendix             % inserts the methods appendix
%
% The figure directory can be overridden before this file is input:
%   \renewcommand{\cnnbgofigdir}{path/to/cnn_bgo_win}

\providecommand{\cnnbgofigdir}{../analysis/results/cnn/cnn_bgo_win}
\providecommand{\cnncode}[1]{\texttt{#1}}
\providecommand{\cnnpath}[1]{\path{#1}}

% =====================================================================
% Main-text section

\section{Position resolution in cooled BGO with a convolutional network}
\label{sec:cnn-bgo}

The three-dimensional interaction position was reconstructed from the
$8\times8$ SiPM light map with CryspNet, a compact convolutional neural network.
The result reported here uses the complete held-out BGO sample in the standard
511-keV analysis window. ``No selection'' means that every event in this window
is retained: no reconstructability or CNN-quality requirement is applied. The
test sample contains $54\,021$ events and includes the photoelectric and Compton
topologies present in the window.

The residual distributions have a narrow central population and a broader tail,
so a single RMS does not fully characterize the response. The central 68\%
absolute-residual quantiles are $1.08$, $1.08$ and $1.26$~mm in $x$, $y$ and
$z$, respectively; the corresponding global RMS values are $1.62$, $1.62$ and
$2.46$~mm. In three dimensions, $9.73\%$ of the reconstructed points lie more
than 5~mm from the true first interaction.

A compact detector-response description is obtained by fitting each coordinate
with two Gaussian components sharing a mean,
\begin{equation}
  p(\Delta q) = f_{\rm c}\,
  \mathcal{N}(\mu,\sigma_{\rm c})
  + (1-f_{\rm c})\,
  \mathcal{N}(\mu,\sigma_{\rm t}),
  \qquad q\in\{x,y,z\}.
  \label{eq:cnn-bgo-two-gauss}
\end{equation}
Table~\ref{tab:cnn-bgo-resolution} gives the complete no-selection result. The
Gaussian core has widths $0.80$, $0.82$ and $0.96$~mm, corresponding to FWHM
values $1.88$, $1.92$ and $2.26$~mm. It contains $72$--$82\%$ of the events,
while the remaining response is described by a component with
$\sigma_{\rm t}\simeq2.5$--$2.6$~mm. Thus cooled BGO provides a sub-millimetre
Gaussian core in all three coordinates and a typical-event resolution of about
1.1--1.3~mm when the full, unselected window sample is considered.

\begin{table}[htbp]
  \centering
  \begin{tabular}{lrrrrrr}
    \toprule
    coordinate
      & $p_{68}$ & RMS
      & $\sigma_{\rm c}$ & FWHM$_{\rm c}$
      & $f_{\rm c}$ & $\sigma_{\rm t}$ \\
      & \multicolumn{2}{c}{[mm]}
      & \multicolumn{2}{c}{[mm]}
      & & [mm] \\
    \midrule
    $x$ & 1.083 & 1.619 & 0.796 & 1.875 & 0.716 & 2.505 \\
    $y$ & 1.081 & 1.622 & 0.816 & 1.922 & 0.737 & 2.586 \\
    $z$ & 1.256 & 2.464 & 0.959 & 2.258 & 0.821 & 2.588 \\
    \bottomrule
  \end{tabular}
  \caption{No-selection CNN position response in BGO. The first two columns
  describe the full residual distribution; the remaining columns are the
  two-Gaussian parameters of Eq.~\eqref{eq:cnn-bgo-two-gauss}.}
  \label{tab:cnn-bgo-resolution}
\end{table}

\begin{figure}[htbp]
  \centering
  \includegraphics[width=\textwidth]{\cnnbgofigdir/residual_fit.png}
  \caption{CNN residuals for the complete no-selection BGO test sample in
  $x$, $y$ and $z$ (steps). The solid curve is the two-component Gaussian fit;
  the dashed and dotted curves show its core and tail terms. The logarithmic
  ordinate exposes the broad response beneath the sub-millimetre core. The
  modest positive skew in the depth residual is retained in the data and is
  summarized, rather than reproduced in detail, by the symmetric fit.}
  \label{fig:cnn-bgo-residual-fit}
\end{figure}

% =====================================================================
% Appendix block. Invoke with \cnnbgoappendix after the parent paper's
% \appendix command.

\newcommand{\cnnbgoappendix}{%
\section{CNN reconstruction and response fit}
\label{app:cnn-bgo}

\subsection{Network and inputs}

Each event is represented by its $8\times8$ photoelectron map and its total
photoelectron count. The map is divided by its sum, separating the spatial light
pattern from the event amplitude. The standardized logarithm of the total count
is supplied as a second, scalar input. The three target coordinates are mapped
linearly to $[-1,1]$ during training and converted back to millimetres for the
reported residuals.

CryspNet begins with a $3\times3$ convolution producing 64 channels, followed by
two 64-channel residual blocks at the native $8\times8$ resolution. A stride-two
convolution increases the representation to 128 channels and reduces it to
$4\times4$; two further residual blocks are followed by global average pooling.
The pooled 128-component descriptor is concatenated with the scalar light-yield
input and passed through a 128-unit fully connected layer and a linear
three-coordinate output. Batch normalization and ReLU activations are used
throughout the convolutional path. The network contains $830\,531$ trainable
parameters.

\subsection{Training outcome}

The window sample contains $360\,135$ events. A seeded permutation divides it
into 70\% training, 15\% validation and 15\% test subsets; the test subset is
held out until the final evaluation. The square detector symmetry supplies the
eight exact rotations and reflections of every training map, with the same
transformations applied to the transverse target coordinates. Validation and
test events remain unaugmented.

Training minimizes a Huber loss on the normalized coordinates with AdamW,
using a batch size of 512, an initial learning rate of $10^{-3}$, weight decay
$10^{-4}$ and a cosine learning-rate schedule over 15 epochs. The training loss
falls from $0.0500$ to $0.0390$. The saved checkpoint is the epoch-10 state,
which minimizes the mean validation RMSE; its coordinate validation values are
$1.608$, $1.612$ and $2.434$~mm. The independent test result is the response
reported in Table~\ref{tab:cnn-bgo-resolution}.

\subsection{Two-component response fit}

The fit uses the CNN prediction minus the true first-interaction coordinate for
each of the $54\,021$ test events. Each residual is histogrammed in 150 bins over
$[-15,15]$~mm and fitted with Eq.~\eqref{eq:cnn-bgo-two-gauss}, using
Poisson weights. The two components share the fitted mean; the narrower width is
identified as the core. The fitted means are $-0.032$, $-0.058$ and
$-0.127$~mm for $x$, $y$ and $z$, showing a small global bias. The depth
distribution also contains a positive, non-Gaussian shoulder from multi-site
events whose light is dominated by a deeper interaction. Equation
\eqref{eq:cnn-bgo-two-gauss} supplies a compact core--tail parameterization,
while Fig.~\ref{fig:cnn-bgo-residual-fit} displays the full residual shape.

The stored configuration is
\cnnpath{analysis/cnn/runs/cnn_bgo_win.toml}. The training history, metrics,
held-out events and fit products are stored under
\cnnpath{analysis/results/cnn/cnn_bgo_win/} as \cnnpath{history.json},
\cnnpath{metrics.json}, \cnnpath{test_events.h5}, \cnnpath{residual_fit.json}
and \cnnpath{residual_fit.png}. The response fit is produced from the held-out
event table by \cnnpath{analysis/cnn/fit_residuals.py}; it does not retrain the
network.
}
             % inserts the main-text section below
%   ...
%   \appendix
%   \cnnbgoappendix             % inserts the methods appendix
%
% The figure directory can be overridden before this file is input:
%   \renewcommand{\cnnbgofigdir}{path/to/cnn_bgo_win}

\providecommand{\cnnbgofigdir}{../analysis/results/cnn/cnn_bgo_win}
\providecommand{\cnncode}[1]{\texttt{#1}}
\providecommand{\cnnpath}[1]{\path{#1}}

% =====================================================================
% Main-text section

\section{Position resolution in cooled BGO with a convolutional network}
\label{sec:cnn-bgo}

The three-dimensional interaction position was reconstructed from the
$8\times8$ SiPM light map with CryspNet, a compact convolutional neural network.
The result reported here uses the complete held-out BGO sample in the standard
511-keV analysis window. ``No selection'' means that every event in this window
is retained: no reconstructability or CNN-quality requirement is applied. The
test sample contains $54\,021$ events and includes the photoelectric and Compton
topologies present in the window.

The residual distributions have a narrow central population and a broader tail,
so a single RMS does not fully characterize the response. The central 68\%
absolute-residual quantiles are $1.08$, $1.08$ and $1.26$~mm in $x$, $y$ and
$z$, respectively; the corresponding global RMS values are $1.62$, $1.62$ and
$2.46$~mm. In three dimensions, $9.73\%$ of the reconstructed points lie more
than 5~mm from the true first interaction.

A compact detector-response description is obtained by fitting each coordinate
with two Gaussian components sharing a mean,
\begin{equation}
  p(\Delta q) = f_{\rm c}\,
  \mathcal{N}(\mu,\sigma_{\rm c})
  + (1-f_{\rm c})\,
  \mathcal{N}(\mu,\sigma_{\rm t}),
  \qquad q\in\{x,y,z\}.
  \label{eq:cnn-bgo-two-gauss}
\end{equation}
Table~\ref{tab:cnn-bgo-resolution} gives the complete no-selection result. The
Gaussian core has widths $0.80$, $0.82$ and $0.96$~mm, corresponding to FWHM
values $1.88$, $1.92$ and $2.26$~mm. It contains $72$--$82\%$ of the events,
while the remaining response is described by a component with
$\sigma_{\rm t}\simeq2.5$--$2.6$~mm. Thus cooled BGO provides a sub-millimetre
Gaussian core in all three coordinates and a typical-event resolution of about
1.1--1.3~mm when the full, unselected window sample is considered.

\begin{table}[htbp]
  \centering
  \begin{tabular}{lrrrrrr}
    \toprule
    coordinate
      & $p_{68}$ & RMS
      & $\sigma_{\rm c}$ & FWHM$_{\rm c}$
      & $f_{\rm c}$ & $\sigma_{\rm t}$ \\
      & \multicolumn{2}{c}{[mm]}
      & \multicolumn{2}{c}{[mm]}
      & & [mm] \\
    \midrule
    $x$ & 1.083 & 1.619 & 0.796 & 1.875 & 0.716 & 2.505 \\
    $y$ & 1.081 & 1.622 & 0.816 & 1.922 & 0.737 & 2.586 \\
    $z$ & 1.256 & 2.464 & 0.959 & 2.258 & 0.821 & 2.588 \\
    \bottomrule
  \end{tabular}
  \caption{No-selection CNN position response in BGO. The first two columns
  describe the full residual distribution; the remaining columns are the
  two-Gaussian parameters of Eq.~\eqref{eq:cnn-bgo-two-gauss}.}
  \label{tab:cnn-bgo-resolution}
\end{table}

\begin{figure}[htbp]
  \centering
  \includegraphics[width=\textwidth]{\cnnbgofigdir/residual_fit.png}
  \caption{CNN residuals for the complete no-selection BGO test sample in
  $x$, $y$ and $z$ (steps). The solid curve is the two-component Gaussian fit;
  the dashed and dotted curves show its core and tail terms. The logarithmic
  ordinate exposes the broad response beneath the sub-millimetre core. The
  modest positive skew in the depth residual is retained in the data and is
  summarized, rather than reproduced in detail, by the symmetric fit.}
  \label{fig:cnn-bgo-residual-fit}
\end{figure}

% =====================================================================
% Appendix block. Invoke with \cnnbgoappendix after the parent paper's
% \appendix command.

\newcommand{\cnnbgoappendix}{%
\section{CNN reconstruction and response fit}
\label{app:cnn-bgo}

\subsection{Network and inputs}

Each event is represented by its $8\times8$ photoelectron map and its total
photoelectron count. The map is divided by its sum, separating the spatial light
pattern from the event amplitude. The standardized logarithm of the total count
is supplied as a second, scalar input. The three target coordinates are mapped
linearly to $[-1,1]$ during training and converted back to millimetres for the
reported residuals.

CryspNet begins with a $3\times3$ convolution producing 64 channels, followed by
two 64-channel residual blocks at the native $8\times8$ resolution. A stride-two
convolution increases the representation to 128 channels and reduces it to
$4\times4$; two further residual blocks are followed by global average pooling.
The pooled 128-component descriptor is concatenated with the scalar light-yield
input and passed through a 128-unit fully connected layer and a linear
three-coordinate output. Batch normalization and ReLU activations are used
throughout the convolutional path. The network contains $830\,531$ trainable
parameters.

\subsection{Training outcome}

The window sample contains $360\,135$ events. A seeded permutation divides it
into 70\% training, 15\% validation and 15\% test subsets; the test subset is
held out until the final evaluation. The square detector symmetry supplies the
eight exact rotations and reflections of every training map, with the same
transformations applied to the transverse target coordinates. Validation and
test events remain unaugmented.

Training minimizes a Huber loss on the normalized coordinates with AdamW,
using a batch size of 512, an initial learning rate of $10^{-3}$, weight decay
$10^{-4}$ and a cosine learning-rate schedule over 15 epochs. The training loss
falls from $0.0500$ to $0.0390$. The saved checkpoint is the epoch-10 state,
which minimizes the mean validation RMSE; its coordinate validation values are
$1.608$, $1.612$ and $2.434$~mm. The independent test result is the response
reported in Table~\ref{tab:cnn-bgo-resolution}.

\subsection{Two-component response fit}

The fit uses the CNN prediction minus the true first-interaction coordinate for
each of the $54\,021$ test events. Each residual is histogrammed in 150 bins over
$[-15,15]$~mm and fitted with Eq.~\eqref{eq:cnn-bgo-two-gauss}, using
Poisson weights. The two components share the fitted mean; the narrower width is
identified as the core. The fitted means are $-0.032$, $-0.058$ and
$-0.127$~mm for $x$, $y$ and $z$, showing a small global bias. The depth
distribution also contains a positive, non-Gaussian shoulder from multi-site
events whose light is dominated by a deeper interaction. Equation
\eqref{eq:cnn-bgo-two-gauss} supplies a compact core--tail parameterization,
while Fig.~\ref{fig:cnn-bgo-residual-fit} displays the full residual shape.

The stored configuration is
\cnnpath{analysis/cnn/runs/cnn_bgo_win.toml}. The training history, metrics,
held-out events and fit products are stored under
\cnnpath{analysis/results/cnn/cnn_bgo_win/} as \cnnpath{history.json},
\cnnpath{metrics.json}, \cnnpath{test_events.h5}, \cnnpath{residual_fit.json}
and \cnnpath{residual_fit.png}. The response fit is produced from the held-out
event table by \cnnpath{analysis/cnn/fit_residuals.py}; it does not retrain the
network.
}
             % inserts the main-text section below
%   ...
%   \appendix
%   \cnnbgoappendix             % inserts the methods appendix
%
% The figure directory can be overridden before this file is input:
%   \renewcommand{\cnnbgofigdir}{path/to/cnn_bgo_win}

\providecommand{\cnnbgofigdir}{../analysis/results/cnn/cnn_bgo_win}
\providecommand{\cnncode}[1]{\texttt{#1}}
\providecommand{\cnnpath}[1]{\path{#1}}

% =====================================================================
% Main-text section

\section{Position resolution in cooled BGO with a convolutional network}
\label{sec:cnn-bgo}

The three-dimensional interaction position was reconstructed from the
$8\times8$ SiPM light map with CryspNet, a compact convolutional neural network.
The result reported here uses the complete held-out BGO sample in the standard
511-keV analysis window. ``No selection'' means that every event in this window
is retained: no reconstructability or CNN-quality requirement is applied. The
test sample contains $54\,021$ events and includes the photoelectric and Compton
topologies present in the window.

The residual distributions have a narrow central population and a broader tail,
so a single RMS does not fully characterize the response. The central 68\%
absolute-residual quantiles are $1.08$, $1.08$ and $1.26$~mm in $x$, $y$ and
$z$, respectively; the corresponding global RMS values are $1.62$, $1.62$ and
$2.46$~mm. In three dimensions, $9.73\%$ of the reconstructed points lie more
than 5~mm from the true first interaction.

A compact detector-response description is obtained by fitting each coordinate
with two Gaussian components sharing a mean,
\begin{equation}
  p(\Delta q) = f_{\rm c}\,
  \mathcal{N}(\mu,\sigma_{\rm c})
  + (1-f_{\rm c})\,
  \mathcal{N}(\mu,\sigma_{\rm t}),
  \qquad q\in\{x,y,z\}.
  \label{eq:cnn-bgo-two-gauss}
\end{equation}
Table~\ref{tab:cnn-bgo-resolution} gives the complete no-selection result. The
Gaussian core has widths $0.80$, $0.82$ and $0.96$~mm, corresponding to FWHM
values $1.88$, $1.92$ and $2.26$~mm. It contains $72$--$82\%$ of the events,
while the remaining response is described by a component with
$\sigma_{\rm t}\simeq2.5$--$2.6$~mm. Thus cooled BGO provides a sub-millimetre
Gaussian core in all three coordinates and a typical-event resolution of about
1.1--1.3~mm when the full, unselected window sample is considered.

\begin{table}[htbp]
  \centering
  \begin{tabular}{lrrrrrr}
    \toprule
    coordinate
      & $p_{68}$ & RMS
      & $\sigma_{\rm c}$ & FWHM$_{\rm c}$
      & $f_{\rm c}$ & $\sigma_{\rm t}$ \\
      & \multicolumn{2}{c}{[mm]}
      & \multicolumn{2}{c}{[mm]}
      & & [mm] \\
    \midrule
    $x$ & 1.083 & 1.619 & 0.796 & 1.875 & 0.716 & 2.505 \\
    $y$ & 1.081 & 1.622 & 0.816 & 1.922 & 0.737 & 2.586 \\
    $z$ & 1.256 & 2.464 & 0.959 & 2.258 & 0.821 & 2.588 \\
    \bottomrule
  \end{tabular}
  \caption{No-selection CNN position response in BGO. The first two columns
  describe the full residual distribution; the remaining columns are the
  two-Gaussian parameters of Eq.~\eqref{eq:cnn-bgo-two-gauss}.}
  \label{tab:cnn-bgo-resolution}
\end{table}

\begin{figure}[htbp]
  \centering
  \includegraphics[width=\textwidth]{\cnnbgofigdir/residual_fit.png}
  \caption{CNN residuals for the complete no-selection BGO test sample in
  $x$, $y$ and $z$ (steps). The solid curve is the two-component Gaussian fit;
  the dashed and dotted curves show its core and tail terms. The logarithmic
  ordinate exposes the broad response beneath the sub-millimetre core. The
  modest positive skew in the depth residual is retained in the data and is
  summarized, rather than reproduced in detail, by the symmetric fit.}
  \label{fig:cnn-bgo-residual-fit}
\end{figure}

% =====================================================================
% Appendix block. Invoke with \cnnbgoappendix after the parent paper's
% \appendix command.

\newcommand{\cnnbgoappendix}{%
\section{CNN reconstruction and response fit}
\label{app:cnn-bgo}

\subsection{Network and inputs}

Each event is represented by its $8\times8$ photoelectron map and its total
photoelectron count. The map is divided by its sum, separating the spatial light
pattern from the event amplitude. The standardized logarithm of the total count
is supplied as a second, scalar input. The three target coordinates are mapped
linearly to $[-1,1]$ during training and converted back to millimetres for the
reported residuals.

CryspNet begins with a $3\times3$ convolution producing 64 channels, followed by
two 64-channel residual blocks at the native $8\times8$ resolution. A stride-two
convolution increases the representation to 128 channels and reduces it to
$4\times4$; two further residual blocks are followed by global average pooling.
The pooled 128-component descriptor is concatenated with the scalar light-yield
input and passed through a 128-unit fully connected layer and a linear
three-coordinate output. Batch normalization and ReLU activations are used
throughout the convolutional path. The network contains $830\,531$ trainable
parameters.

\subsection{Training outcome}

The window sample contains $360\,135$ events. A seeded permutation divides it
into 70\% training, 15\% validation and 15\% test subsets; the test subset is
held out until the final evaluation. The square detector symmetry supplies the
eight exact rotations and reflections of every training map, with the same
transformations applied to the transverse target coordinates. Validation and
test events remain unaugmented.

Training minimizes a Huber loss on the normalized coordinates with AdamW,
using a batch size of 512, an initial learning rate of $10^{-3}$, weight decay
$10^{-4}$ and a cosine learning-rate schedule over 15 epochs. The training loss
falls from $0.0500$ to $0.0390$. The saved checkpoint is the epoch-10 state,
which minimizes the mean validation RMSE; its coordinate validation values are
$1.608$, $1.612$ and $2.434$~mm. The independent test result is the response
reported in Table~\ref{tab:cnn-bgo-resolution}.

\subsection{Two-component response fit}

The fit uses the CNN prediction minus the true first-interaction coordinate for
each of the $54\,021$ test events. Each residual is histogrammed in 150 bins over
$[-15,15]$~mm and fitted with Eq.~\eqref{eq:cnn-bgo-two-gauss}, using
Poisson weights. The two components share the fitted mean; the narrower width is
identified as the core. The fitted means are $-0.032$, $-0.058$ and
$-0.127$~mm for $x$, $y$ and $z$, showing a small global bias. The depth
distribution also contains a positive, non-Gaussian shoulder from multi-site
events whose light is dominated by a deeper interaction. Equation
\eqref{eq:cnn-bgo-two-gauss} supplies a compact core--tail parameterization,
while Fig.~\ref{fig:cnn-bgo-residual-fit} displays the full residual shape.

The stored configuration is
\cnnpath{analysis/cnn/runs/cnn_bgo_win.toml}. The training history, metrics,
held-out events and fit products are stored under
\cnnpath{analysis/results/cnn/cnn_bgo_win/} as \cnnpath{history.json},
\cnnpath{metrics.json}, \cnnpath{test_events.h5}, \cnnpath{residual_fit.json}
and \cnnpath{residual_fit.png}. The response fit is produced from the held-out
event table by \cnnpath{analysis/cnn/fit_residuals.py}; it does not retrain the
network.
}
